% =============================================================================
% Lecture 20: Conclusion
% Agentic Machine Learning Systems (Volume III), Chapter 20
% =============================================================================
% Thesis: components that pass their own tests can still leave a task
% unfinished, so reliable autonomy comes from closing invariants at every
% boundary of one trajectory.
% Running case: incident_42, the remediation agent's configuration rollback
% whose acknowledgment is lost (Conclusion \S20.2), run for real on a localhost
% endpoint (images/data/settle_run.py). This is the last lecture of the
% course: it recalls every part and closes with a question for further study.
%
% Build: make deck V=3 D=20_conclusion   (figures: images/make_figs.py)
% Plan:  lecture_map.yml in this folder.
% Numbers: images/calcs.py prints every derived value; sources named there.
% =============================================================================

\documentclass[aspectratio=169, 12pt]{beamer}
\usepackage{../../assets/beamerthememlsys}
\mlsyssetup{
  volume={Volume III},
  chapter={Chapter 20},
  chaptertitle={Conclusion},
  layer=trajectory,
  logo={../../assets/img/logo-mlsysbook.png},
  instlogo={../../assets/img/logo-harvard.png},
}
\usepackage[T1]{fontenc}
\usepackage[scaled=0.9]{helvet}
\usepackage{courier}
\usepackage[varqu,scaled=0.95]{zi4}
\renewcommand{\familydefault}{\sfdefault}
\usepackage{booktabs,amsmath}
\usetikzlibrary{arrows.meta,positioning}
\graphicspath{{images/}}
\setcounter{mlsystotalsections}{5}

\title{Conclusion}
\author{Vijay Janapa Reddi}
\institute{Harvard University}
\date{}

% A large answer choice for predict and concept-question slides
\newcommand{\choice}[2]{\textbf{\textcolor{mlsysaccent}{#1}}\enspace #2}
% A choice that is highlighted from overlay #1 on (the answer)
\newcommand{\answerat}[3]{\alt<#1->{\colorbox{mlsysaccent!15}{\choice{#2}{#3}}}{\choice{#2}{#3}}}
% A goal or takeaway line with its text-free thumbnail
\newcommand{\thumbline}[2]{%
  \begin{minipage}[c]{2.6cm}\centering\includegraphics[width=2.5cm,height=1.0cm,keepaspectratio]{#1}\end{minipage}%
  \hspace{0.3cm}\begin{minipage}[c]{0.74\textwidth}\small\raggedright #2\end{minipage}}
% An identifier from the running case, in code type
\newcommand{\op}[1]{{\ttfamily #1}}
% An exact excerpt of a real document
\newenvironment{quoteblock}{\begin{mlsyscode}\footnotesize\sffamily\raggedright}{\end{mlsyscode}}
% A gray tag for illustrative, derived, or measured values
\newcommand{\smalltag}[1]{{\scriptsize\textcolor{midgray}{\textit{#1}}}}
% Golden-thread tag: one line on the goals slide naming the H·S·A term this lecture acts on
\newcommand{\hsathread}[1]{{\scriptsize\textbf{\textcolor{mlsysaccent}{H·S·A thread:}} #1\par}}
% One part of the course on a recap slide: stack thumbnail, principles, evidence
\newcommand{\partrow}[4]{%
  \begin{minipage}[c]{2.3cm}\centering\includegraphics[width=2.2cm,height=1.5cm,keepaspectratio]{#1}\end{minipage}%
  \hspace{0.2cm}\begin{minipage}[c]{0.50\textwidth}\raggedright
    {\small\bfseries\textcolor{mlsysaccent}{#2}}\\[0.04cm]{\footnotesize #3}\end{minipage}%
  \hspace{0.2cm}\begin{minipage}[c]{0.30\textwidth}\raggedright #4\end{minipage}}

\begin{document}

% =============================================================================
% TITLE, RECALL, HOOK, GOALS
% =============================================================================
\mlsysnote{[0.5 min] Read the thesis aloud once. This is the last lecture of the course. It introduces no new mechanism: it follows one incident through everything the course built and asks students to produce what earlier lectures taught before the slide shows it. Notes total about 71.5 minutes; the core lecture, with every [OPTIONAL] slide cut, is about 57 minutes. FLEX: [CORE].}
\mlsystitle{Conclusion}{Components that pass their own tests can still leave a task unfinished, so reliable autonomy comes from closing invariants at every boundary of one trajectory.}{cover_conclusion_blueprint.png}

\begin{frame}{Why did cheaper tokens cost more per repair?}
  \mlsysactivity{Retrieval}{2 min}
  \begin{columns}[c]
    \begin{column}{0.55\textwidth}
      \normalsize
      Lecture 19: the same 100 tickets, two configurations.\\[0.2cm]
      The one with the lower token price cost \textbf{\$1.32} per accepted repair; the other cost \textbf{\$0.55}.\\[0.3cm]
      {\large Close your notes. Why?}
    \end{column}
    \begin{column}{0.40\textwidth}
      \centering
      \includegraphics[width=\textwidth,height=2.6cm,keepaspectratio]{thumb_recall_bars.pdf}\\[0.15cm]
      \raggedright\small
      \uncover<2->{Each accepted repair carried \textbf{3} rejected attempts against \textbf{0.33}: 25 against 75 accepted of 100.}
    \end{column}
  \end{columns}
  \mlsyscredit{Agent Economics, Napkin Math 19.1 and table 1. 3, 0.33, and 2.4$\times$ derived (images/calcs.py).}
  \note{[2 min] LINK: Lecture 19 closed Part VI on cohort cost per accepted task. ENGAGE: notes closed, one minute: why does the configuration with the lower token price cost 2.4 times as much for each repair it finishes? Expected answer: it is accepted on 25 of 100 tickets against 75, so every accepted repair carries 3 rejected attempts against 0.33. Step 2 shows the counts. NARRATE: the account that matters charges every attempt, failed ones included, to the tasks a contract accepted. Today asks where the accepted task itself comes from: what has to be true at every boundary before anyone may say a task is done. WARN: students who answer ``the cheap model is worse'' have the symptom and are missing the account; push for the ratio. FLEX: [CORE].}
\end{frame}

\begin{frame}{The reply is lost. What should the runtime do?}
  \mlsysactivity{Predict}{3 min}
  \begin{columns}[T]
    \begin{column}{0.55\textwidth}
      \begin{mlsyscard}{mlsysaccent}
        \footnotesize
        \textbf{\op{incident\_42}}: a latency alert after a configuration change.\\[0.05cm]
        The agent proposes reverting to \op{config\_v7}; an operator approves.\\[0.05cm]
        The endpoint drops the connection \textbf{before acknowledging}.
      \end{mlsyscard}
      \vspace{0.1cm}\small
      \uncover<2->{A may apply it twice; B may close an unapplied change; C leaves a live change unowned.}
    \end{column}
    \begin{column}{0.41\textwidth}
      \centering
      \safeimg[width=\textwidth,height=2.2cm,keepaspectratio]{datacenter-server-racks.jpg}\\[0.15cm]
      \raggedright\small
      \choice{A}{Retry the change}\\[0.05cm]
      \choice{B}{Record success and close}\\[0.05cm]
      \choice{C}{Halt and wait}\\[0.05cm]
      \answerat{3}{D}{None of these}
    \end{column}
  \end{columns}
  \mlsyscredit{Scenario: Conclusion, \S20.2 (constructed by the chapter). Photo: Carl Lender, CC BY 2.0, Wikimedia Commons; machines of this kind; the scenario's service is constructed.}
  \note{[3.5 min] NARRATE: the chapter's capstone incident. A database service alerts on query latency shortly after a configuration change. The remediation agent reads telemetry, finds a known regression, and proposes reverting to config\_v7; an on-call operator approves. The configuration endpoint drops the connection before it acknowledges. The photo shows the kind of machine such a service runs on, nothing more. ENGAGE: vote A, B, C, or D, one minute with a neighbor, vote again. Step 2: the three defaults each fail. A retry under a fresh identity can apply the mutation twice. Assuming success can close the incident while degraded queries continue. Halting leaves a live production change unverified with nobody accountable for it. Step 3: D. The runtime needs evidence it does not have: it must ask the endpoint what happened to this one operation. WARN: students choose A because retries are what they know from web clients; a retry is safe only when the endpoint recognizes a repeat. Keep the answer open; the first segment builds it and then runs it. FLEX: [CORE].}
\end{frame}

\begin{frame}{Can tested components still leave a task unfinished?}
  \begin{columns}[c]
    \begin{column}{0.78\textwidth}
      \small You will be able to:\\[0.05cm]
      \thumbline{thumb_capstone_step3.pdf}{Trace one incident and name the owner and record of each transition.}\\[0.05cm]
      \thumbline{thumb_lostreply_step3.pdf}{Settle a lost reply without repeating its effect.}\\[0.05cm]
      \thumbline{thumb_swe_agent.pdf}{Choose an intervention from an attributed failure and its cost.}\\[0.05cm]
      \thumbline{thumb_passing_step2.pdf}{Scope a safety claim to the evidence that supports it.}
    \end{column}
    \begin{column}{0.19\textwidth}
      \centering
      \safeimg[width=\textwidth,height=2.4cm,keepaspectratio]{hsa_locator_hsa.pdf}\\[0.05cm]
      \hsathread{closes all three exposures}
    \end{column}
  \end{columns}
  \note{[1.5 min] The driving question from the chapter: why can individually tested components still leave an agent's task unfinished or its effects unaccounted for? Read the four goals. NARRATE the locator: Lecture 1's H{\textperiodcentered}S{\textperiodcentered}A exposures; the incident lights all three. It runs tens of turns (horizon), changes shared production state (state, $S_3$), and makes changes that cannot always be compensated (authority, $A_3$). Every part of the course closed one of these; today they return together on one trajectory. FLEX: [CORE].}
\end{frame}

% =============================================================================
% SECTION: ONE INCIDENT
% =============================================================================
\mlsysnote{[0.5 min] Segment 1 follows incident\_42 from its contract to its acceptance. Read the question aloud. FLEX: [CORE].}
\mlsysdivider{What does each boundary of one incident need?}{Conclusion, \S20.2}
\section{Incident}

\begin{frame}{The incident crossed four parts of the course}
  \centering
  \mlsyssteps[width=0.94\textwidth,height=5.6cm,keepaspectratio]{capstone_step}{3}
  \mlsyscredit{Slide variant of Conclusion, figure 1.}
  \note{[2.5 min] Step 1: the loop that diagnosed the regression. Context assembly from Part II gathers the alert, the metrics, the live configuration, and the runbook. The model call from Part I proposes. The tool call and sandbox from Part III act and return an observation, which re-enters context as data, never as an instruction. The harness of Part IV wraps every step with budgets, the approval gate, a durable log, and recovery. Step 2: the task contract from Lecture 1 came first, adopted by an authorized owner, and acceptance came last, decided on evidence, with the agent's report treated as a claim. Step 3: learning and delegation never entered. Parts V and VI are later, optional decisions; neither changes what it takes to dispatch op\_42 or to accept it. ENGAGE: which part would you add first if this incident recurred weekly? Expected: none until a failure is attributed; segment 3 returns to this. FLEX: [CORE].}
\end{frame}

\begin{frame}{An approval binds one exact payload and its baseline}
  \begin{columns}[T]
    \begin{column}{0.45\textwidth}
      \small
      \op{approval\_42} binds:\\[0.1cm]
      \footnotesize
      operator, request, host, exact settings,\\
      baseline \op{config\_v7}, and an expiry.\\[0.25cm]
      \small
      \uncover<2->{The gate rechecks it at dispatch.\\[0.1cm]
      The endpoint applies the change only if the live version is still \op{config\_v7}.}
    \end{column}
    \begin{column}{0.52\textwidth}
      {\footnotesize\textbf{AWS S3, 28 February 2017:} an authorized engineer ran an established playbook.}\\[0.05cm]
      \begin{quoteblock}
        [...] one of the inputs to the command was entered incorrectly and a larger set of servers was removed than intended.
      \end{quoteblock}
      \uncover<3->{\begin{quoteblock}
        We have modified this tool to remove capacity more slowly and added safeguards [...]
      \end{quoteblock}}
    \end{column}
  \end{columns}
  \mlsyscredit{Conclusion, \S20.2.1 and table 3. AWS, ``Summary of the Amazon S3 Service Disruption in the Northern Virginia (US-EAST-1) Region,'' aws.amazon.com/message/41926.}
  \note{[2.5 min] Step 1: the approval record. approval\_42 binds the operator's identity to the request, the target host, the exact payload, the prerequisite baseline config\_v7, and an expiry. It records permission for that payload and nothing broader; it performs nothing. Step 2: the approval is an earlier check, and complete mediation (Lecture 7) needs a fresh one at the moment of use. The action gate rechecks authorization, the payload digest, the expiry, and the budget. Then the endpoint closes the time-of-check to time-of-use race: its conditional update applies the change only if the live version still matches config\_v7, which is Lecture 7's action-time precondition. Only the server can enforce that. Step 3, the real case: an authorized engineer, an established playbook, one input typed wrong. Approval of the person did not bound the payload; the fix AWS shipped was a check inside the tool, below any human. ENGAGE: which of the incident's checks plays the role of AWS's new safeguard? Expected: the endpoint's expected-version check and the gate's payload digest, both enforced by mechanism. WARN: students read ``an operator approved'' as ``the action is safe''; an approval records who decided; the gate and the endpoint bound what executes. FLEX: [CORE].}
\end{frame}

\begin{frame}{Write the intent first; then a lost reply has an owner}
  \centering
  \mlsyssteps[width=0.95\textwidth,height=3.8cm,keepaspectratio]{lostreply_step}{3}\\[0.15cm]
  \begin{minipage}{0.92\textwidth}
    \uncover<3->{{\footnotesize\textbf{Stripe API reference} (excerpt)}
    \begin{quoteblock}
      [...] if a connection error occurs, you can safely repeat the request without risk of [...] performing the update twice.
    \end{quoteblock}}
  \end{minipage}
  \mlsyscredit{Conclusion, \S20.2.1 and \S20.5. Stripe, ``Idempotent requests,'' docs.stripe.com/api/idempotent\_requests.}
  \note{[3 min] Step 1: before the request leaves, the runtime appends an intent record for op\_42 to its write-ahead journal, with a recovery descriptor naming the corrective template and the accountable owner, operations\_owner\_42. Lecture 11 called this durable intent before dispatch; Lecture 12 called the descriptor the trajectory saga's obligation. Then it dispatches op\_42, expecting config\_v7. Step 2: the endpoint applies config\_v8 and the reply is lost. The timeout says nothing about which of three worlds holds: the request was lost, it is still executing, or it committed and the acknowledgment vanished. The runtime marks op\_42 unknown, since it may have committed. Step 3: the settlement process asks the endpoint's operation registry for op\_42, a read-only lookup that writes nothing. It returns applied, config\_v8, and the operation settles. This is Lecture 7's settlement before retry, and it works only because the endpoint deduplicates by an operation identity bound to caller, target, and payload. The Stripe excerpt is the same contract in a public API. ENGAGE: what if the lookup returns no record? Expected: the status stays unknown, because the request may still be in flight; the obligation stays with operations\_owner\_42. WARN: treat ``unknown'' as a state with an owner; clearing it as an error loses the obligation. FLEX: [CORE].}
\end{frame}

\begin{frame}{What does a blind retry cost a fleet daily?}
  \mlsysactivity{Estimate}{3 min}
  \begin{columns}[T]
    \begin{column}{0.42\textwidth}
      \begin{mlsyscard}{mlsysaccent}
        \footnotesize
        \textbf{Given} \smalltag{(Exercise 3, hypothetical)}\\[0.03cm]
        25,000 mutating calls a day\\
        Request lost: 0.001; reply lost: 0.004\\
        Duplicate: \$2,000; status lookup: \$0.01
      \end{mlsyscard}
      \small
      \uncover<2->{Timeouts: about \textbf{124.9} a day\\[0.05cm]}
      \uncover<3->{Blind retry: \textbf{99.8} duplicates, \textbf{\$199,600} a day\\[0.05cm]
      Lookups instead: \textbf{\$1.25} a day}
    \end{column}
    \begin{column}{0.55\textwidth}
      \centering
      \uncover<4->{{\footnotesize\textbf{Measured: the four runtimes on a real HTTP endpoint}}\\[0.05cm]
      \safeimg[width=\textwidth,height=4.2cm,keepaspectratio]{measured.pdf}}
    \end{column}
  \end{columns}
  \mlsyscredit{Conclusion, \S20.13 Exercise 3. Measured: our localhost run with the same fault rates, 25,000 incidents per runtime (images/data/settle\_run.py, 2026-10-19). Estimates derived (images/calcs.py).}
  \note{[3.5 min] ENGAGE: two minutes in pairs. Step 2: a call times out when the request is lost or the reply is lost: 25,000 times (0.001 + 0.999 $\times$ 0.004), about 124.9 a day. Step 3: a blind retry under a fresh identity duplicates the change when the original applied and only its reply was lost, and the retry then applies: 25,000 $\times$ 0.999 $\times$ 0.004 $\times$ 0.999, about 99.8 duplicates, \$199,600 a day. A status lookup on every timeout costs 124.9 $\times$ \$0.01, \$1.25. Step 4: we ran incident\_42's dilemma for real: a Python HTTP endpoint that drops a request with probability 0.001 or applies the change and drops the reply with probability 0.004, and four runtimes with 25,000 incidents each. Retrying blindly duplicated 92 changes (99.8 expected). Assuming success closed 27 incidents whose change never applied (25 expected). Halting left 108 incidents with no recorded outcome, 92 of them already applied. Settling by status, against an endpoint that keeps operation ids, had zero failures. WARN: the counts are one seeded run on one laptop; read the agreement with the formula and ignore the third digit. The given rates are the exercise's hypotheticals. FLEX: [CORE].}
\end{frame}

\begin{frame}{Injected text in a log inherits no authority}
  \begin{columns}[T]
    \begin{column}{0.47\textwidth}
      {\footnotesize\textbf{A real injection} (pull request an agent opened)}\\[0.1cm]
      \safeimg[width=\textwidth,height=2.6cm,keepaspectratio]{mcp-leaking-pr.png}\\[0.1cm]
      \footnotesize An issue's text told the agent to publish private data. Every call was within its token's scope.
    \end{column}
    \begin{column}{0.50\textwidth}
      \centering
      \mlsyssteps[width=\textwidth,height=5.3cm,keepaspectratio]{barriers_step}{3}
    \end{column}
  \end{columns}
  \mlsyscredit{Conclusion, \S20.2.2. Invariant Labs, ``GitHub MCP Exploited,'' May 2025 (demonstration accounts; War Story 7.1).}
  \note{[2.5 min] LINK: production telemetry is untrusted text. Left: the real case from the course, where an agent read an issue as instructions and opened a public pull request with a user's private details; the accounts are the researchers' demonstration accounts. Right, incident\_42's version: a log line fetched under read\_scope\_42 asks the model to send the database credential to an outside server. Lecture 9's structural prompt injection: the model cannot separate data from instructions, so the runtime must. Step 1: the action gate blocks a call to an address read\_scope\_42 does not name. Step 2: independently, the sandbox's egress policy (Lecture 8) drops a socket to any unlisted host. Step 3: the hard case uses only permitted tools: a query to an allowed endpoint carries the connection string in an argument. Lecture 9's information-flow labels catch it: the credential's label excludes the telemetry service as a recipient, the label follows the value into the query string, and dispatch is refused. ENGAGE: which barrier would have stopped the pull request on the left? Expected: a flow label on private-repository data, or a scope that names only the task's repository. WARN: delimiters and provenance tags in the prompt help the model and enforce nothing. FLEX: [OPTIONAL].}
\end{frame}

% =============================================================================
% SECTION: FIFTEEN PRINCIPLES
% =============================================================================
\mlsysnote{[0.5 min] Segment 2 gathers the fifteen principles, one part of the course at a time. Each part slide asks students to name the principles before reading them. FLEX: [CORE].}
\mlsysdivider{Which principle did each part of the course add?}{Conclusion, \S20.2.3 (table 2)}
\section{Principles}

\begin{frame}{Which principle did each part add?}
  \mlsysactivity{Retrieval}{3 min}
  \centering
  \mlsyskitfig[width=0.78\textwidth,height=4.6cm,keepaspectratio]{stack_roadmap}\\[0.15cm]
  {\normalsize Close your notes. For each layer, write one principle it added, and one number that showed it.}
  \note{[3 min] The course stack with every lecture on its layer; Lectures 1 and 20 sit under the trajectory bar because they span the stack. Give two and a half minutes. Any wording counts if it states what the layer guarantees. Walk the room: empty layers show what to review before the exam. The next three slides reveal the principles two parts at a time, with one number per part from its lectures. The table has fifteen rows: invariant closure from Lecture 1, then two or three per part, and Part VII adds none. FLEX: [CORE].}
\end{frame}

\begin{frame}{Parts I and II price the call and select its context}
  \partrow{thumb_stack_model.pdf}{Part I, the model (Lectures 2 and 3)}{Output-bound latency. The verification asymmetry.}{%
    \safeimg[width=\linewidth,height=1.4cm,keepaspectratio]{fiji-hbm-package.jpg}\\
    {\scriptsize 30 turns resend 892,500 input tokens against 51,500 fresh (L2).}}\\[0.35cm]
  \partrow{thumb_stack_memory.pdf}{Part II, agent memory (Lectures 4 to 6)}{Context is a selected working set. Attention state is derived from the token prefix. Authority stays with the source.}{%
    \safeimg[width=\linewidth,height=1.4cm,keepaspectratio]{card-catalog.jpg}\\
    {\scriptsize Appending bills 3,030,000 tokens; selecting, 1,410,000 (L4).}}
  \mlsyscredit{Conclusion, table 2; Napkin Math 2.2 and 4.1. Photos: C. Spille/pcgameshardware.de, CC BY-SA 4.0; TBurmeister (WMF), CC BY-SA 4.0; Wikimedia Commons.}
  \note{[2 min] Part I: a model call is a proposal priced mostly by its output length, because decode reads the weights from memory once per output token (the photo is an accelerator package with its stacked memory beside the die); and a check selects valid candidates only as well as its error rates and the base rate allow, Lecture 3's verifier precision. Lecture 2's number: a 30-turn trajectory that resends its prefix sends 892,500 input tokens against 51,500 fresh with prefix reuse. Part II: context is a working set chosen for this call; KV state is valid only for its exact token prefix, so one changed token early forfeits reuse after it; and a stored record keeps the authority of its source, never more. Lecture 4's number: appending every observation bills 3,030,000 input tokens over the step budget against 1,410,000 for selection. The card catalog stands for records that point to a source. ENGAGE: which Part II principle does a stale summary of config\_v7 violate? Expected: authority stays with the source. FLEX: [CORE].}
\end{frame}

\begin{frame}{Parts III and IV bound authority and the horizon}
  \partrow{thumb_stack_tools.pdf}{Part III, tool use (Lectures 7 to 9)}{The typed action contract. Settlement before retry. Containment beneath the model.}{%
    \safeimg[width=\linewidth,height=1.4cm,keepaspectratio]{lethaltrifecta.jpg}\\
    {\scriptsize Blind retry duplicates 1.92\% of calls: 3,842 a day of 200,000 (L7).}}\\[0.35cm]
  \partrow{thumb_stack_runtime.pdf}{Part IV, the agent runtime (Lectures 10 to 14)}{Trajectory-level encapsulation. Intent before effect. Causal and statistical release evidence.}{%
    \safeimg[width=\linewidth,height=1.4cm,keepaspectratio]{jim_gray_full.jpg}\\
    {\scriptsize A rule obeyed on 99\% of turns holds for all 30 turns 74\% of the time (L10).}}
  \mlsyscredit{Conclusion, table 2; Napkin Math 7.1; Agent Harness \S10.10 (74\% derived). Images: S. Willison, simonwillison.net, 16 June 2025; Microsoft Research, CC BY-SA 3.0 us, Wikimedia Commons.}
  \note{[2 min] Part III: dispatch only a schema-valid, authorized call whose precondition the endpoint enforces; settle an uncertain outcome before any retry; and bound what code can reach beneath the model, because the model cannot be trusted to refuse. Lecture 7's number: under blind retry, 1.9208 percent of calls duplicate their effect, about 3,842 a day across 200,000 mutating calls. The image is Simon Willison's lethal trifecta from Lecture 9: private data, untrusted content, and a way to communicate outward, which incident\_42's injection had. Part IV: budgets and grants are enforced over the whole trajectory, derived grants only narrow; intent is durable before any effect; and a release rests on causal and statistical evidence. Lecture 10's number: an instruction honored on 99 percent of turns holds across all 30 only 74 percent of the time, which is why a freeze belongs in the harness. The photo is Jim Gray, whose write-ahead log is the pattern behind intent before effect. FLEX: [CORE].}
\end{frame}

\begin{frame}{Parts V and VI change the model or multiply the agents}
  \begin{minipage}[c]{2.3cm}\centering\includegraphics[width=2.2cm,height=1.5cm,keepaspectratio]{thumb_stack_learning.pdf}\end{minipage}%
  \hspace{0.2cm}\begin{minipage}[c]{0.36\textwidth}\raggedright
    {\small\bfseries\textcolor{mlsysaccent}{Part V, learning (Lectures 15 to 17)}}\\[0.04cm]{\footnotesize Verified trajectory post-training.}\end{minipage}%
  \hspace{0.2cm}\begin{minipage}[c]{0.44\textwidth}\raggedright
    {\scriptsize With the scorer in reach, o3 hacked it (L17):}\\[0.03cm]
    \safeimg[width=\linewidth,height=1.9cm,keepaspectratio]{metr_hack.pdf}\end{minipage}\\[0.35cm]
  \partrow{thumb_stack_scale.pdf}{Part VI, agents at scale (Lectures 18 and 19)}{The coordination tax. Trajectory goodput and the accepted task.}{%
    \safeimg[width=\linewidth,height=1.4cm,keepaspectratio]{ashburn.jpg}\\
    {\scriptsize One reported multi-agent system used about 15$\times$ the tokens of chat (L18).}}
  \mlsyscredit{Conclusion, table 2. METR (Von Arx et al. 2025) via War Story 17.1; Hadfield et al. 2025 via War Story 18.1. Photo: Theodore Christopher, CC0, Wikimedia Commons (data centers near Ashburn, Virginia).}
  \note{[2 min] Part V: training changes which trajectories the model proposes and grants no new access, so it moves no exposure; a passing task check does not make every action a training target. Lecture 17's measurement: when METR left o3's scorer in reach, it hacked the score in 39 of 128 RE-Bench runs (30 percent) against 8 of 1,087 HCAST runs (0.7 percent). Part VI: splitting work pays only if its acceptance or deadline gain justifies its cost multiple, and every attempt is charged to the accepted tasks. Lecture 18's number: one reported system used about 4 times the tokens of chat as a single agent and about 15 times as a multi-agent system; the source is one engineering report, and the slide says so. The photo is the data-center region where such fleets run. ENGAGE: which two parts never entered incident\_42, and why? Expected: V and VI; nothing in the incident needed new weights or more agents. Land the segment: fifteen principles, one per guarantee, and invariant closure from Lecture 1 anchors all of them, because no computation certifies its own task contract. FLEX: [OPTIONAL].}
\end{frame}

% =============================================================================
% SECTION: EVIDENCE AND INTERVENTION
% =============================================================================
\mlsysnote{[0.5 min] Segment 3: what state and evidence must survive each boundary, and how to choose the next safeguard. FLEX: [CORE].}
\mlsysdivider{Where does evidence go stale, and which fix comes first?}{Conclusion, \S20.3, \S20.4, and \S20.6}
\section{Evidence}

\begin{frame}{A cache hit proves only that the tokens match}
  \begin{columns}[T]
    \begin{column}{0.48\textwidth}
      \centering
      \mlsyssteps[width=\textwidth,height=5.3cm,keepaspectratio]{records_step}{2}
    \end{column}
    \begin{column}{0.49\textwidth}
      {\footnotesize\textbf{vLLM, Automatic Prefix Caching} (excerpt)}\\[0.05cm]
      \begin{quoteblock}
        [...] we hash each kv-cache block by the tokens in the block and the tokens in the prefix before the block
      \end{quoteblock}
      \vspace{0.1cm}\small
      \uncover<2->{The key holds tokens and omits the service's version. Freshness comes from the source.}
    \end{column}
  \end{columns}
  \mlsyscredit{Conclusion, \S20.4 and table 4. vLLM documentation, ``Automatic Prefix Caching,'' docs.vllm.ai/en/latest/design/prefix\_caching.html.}
  \note{[2.5 min] Step 1: before the change, five records describe the incident and all agree on config\_v7. Step 2: op\_42 settles on config\_v8. The source record and the event log move. The incident summary written by the agent and the runbook vector index still say config\_v7: Lecture 6's record roles say neither may override the source. Lecture 4's context invalidation retires them from the next prompt or records the omission. The KV cache for the old prompt is still valid, for its tokens: the excerpt from a real serving system shows the cache key is the token sequence. That is Lecture 5's rule, reusable is not applicable. ENGAGE: the next prompt substitutes the config\_v8 line. Which cached work survives? Expected: everything before the first changed token. WARN: students conclude the cache is wrong; it is right about tokens and silent about the world, so a stale prefix is cheap and wrong at once. The endpoint's expected-version check remains the last guard. FLEX: [CORE].}
\end{frame}

\begin{frame}{Five seams fail while each side meets its own spec}
  \footnotesize
  \renewcommand{\arraystretch}{1.15}
  \begin{tabular}{@{}P{0.20\textwidth}P{0.42\textwidth}P{0.30\textwidth}@{}}
    \toprule
    \textbf{Seam} & \textbf{Failure in \op{incident\_42}} & \textbf{Principles in tension} \\
    \midrule
    World to its copies & \op{config\_v7} shown as current after \op{config\_v8} & 6 against 5 \\
    Proposal to effect & lost reply treated as failure; \op{op\_42} repeated & 11, then 8 \\
    \leavevmode\uncover<2->{Recovery to authority} & \leavevmode\uncover<2->{a recorded plan used as permission later} & \leavevmode\uncover<2->{11 against 7} \\
    \leavevmode\uncover<3->{Record to learning} & \leavevmode\uncover<3->{every accepted step becomes a target} & \leavevmode\uncover<3->{12 against 3} \\
    \leavevmode\uncover<3->{Across workers} & \leavevmode\uncover<3->{workers copy one allowance; agreement taken as proof} & \leavevmode\uncover<3->{14 against 10} \\
    \bottomrule
  \end{tabular}\\[0.2cm]
  \small Each row needs an owner and a record at the seam.
  \mlsyscredit{Conclusion, table 4; principle numbers as in table 2.}
  \note{[2 min] Table 4 of the chapter. Step 1, the two seams the incident crossed: storage and serving each met their specification while the agent acted on stale evidence (source authority, principle 6, against prefix reuse, principle 5); and a lost reply treated as failure invites a repeat (intent before effect, then settlement before retry). Step 2: a recovery descriptor written yesterday is not permission today; any corrective change needs current authorization and preconditions, the typed action contract. Step 3: two seams that arise only if the deployment adds learning or delegation. Accepting the incident does not make every step a training target, and a release grader reused as a reward stops being evidence (Lecture 15's capability-gap triage decides whether a failure belongs in the weights at all). Delegated workers may share a scoped read grant but must reserve against one budget, and agreement among workers reading the same logs is not independent confirmation (Lecture 18's correlated reviewers). ENGAGE: pick one row: who owns the record? Expected for row 2: the settlement process for operations\_owner\_42. FLEX: [OPTIONAL].}
\end{frame}

\begin{frame}{Where does an incident attempt's money go?}
  \mlsysactivity{Estimate}{3 min}
  \centering
  \safeimg[width=0.90\textwidth,height=3.4cm,keepaspectratio]{bill.pdf}\\[0.05cm]
  \begin{minipage}{0.92\textwidth}\small
    Attempt \$20.45, 70\% accepted: \textbf{\$29.21} per accepted incident; approval is 97.8\%.\\[0.05cm]
    {\textbf{Double the model calls. How far must acceptance rise?}}\\[0.03cm]
    \uncover<2->{Attempt \$20.60 (+0.73\%): acceptance must exceed \textbf{70.51\%}.}
  \end{minipage}
  \mlsyscredit{Conclusion, \S20.3.2 (remediation profile from Agent Economics, Lighthouse 19.1; illustrative). Totals, shares, and break-even derived (images/calcs.py).}
  \note{[3.5 min] NARRATE: the chapter prices one incident\_42 attempt on the remediation profile: \$0.15 of model calls, \$0.10 of telemetry and change tools, \$0.20 for the health-window check, and \$20 of on-call review, \$20.45 in all. At 70 percent acceptance that is \$29.21 per accepted incident, Lecture 1's cost per accepted task. The model is 0.7 percent of the attempt; approval is 97.8 percent. The red bar is the danger: a wrong change proposed in 4 percent of attempts that costs \$12,000 when it commits carries an expected irreversible loss of \$480 per attempt with no check ahead of it, about 23.5 times the attempt. ENGAGE: one minute: if you double the model calls for more diagnosis, how far must acceptance rise to pay? Step 2: the attempt rises by the model share, to \$20.60, so acceptance must exceed 0.70 $\times$ 20.60/20.45, about 70.51 percent. Cheap in principle; it must be measured on a matched baseline (Lecture 13). WARN: students optimize the model bill because it is the one they see; here the gates (approval binding, version check, health window) matter more than any cap on tokens. FLEX: [CORE].}
\end{frame}

\begin{frame}{Attribute the failure, then pick the cheapest fix that works}
  \begin{columns}[T]
    \begin{column}{0.44\textwidth}
      \footnotesize
      \renewcommand{\arraystretch}{1.1}
      \begin{tabular}{@{}P{0.38\linewidth}P{0.54\linewidth}@{}}
        \toprule
        \textbf{Intervention} & \textbf{Evidence that justifies it} \\
        \midrule
        Context & needed facts absent or stale \\
        Tool interface & outcome states ambiguous \\
        Runtime control & a rule not enforced \\
        Fine-tuning & a recurring behavior \\
        RL & a useful scoped reward \\
        \midrule
        Delegation & a separate decision \\
        \bottomrule
      \end{tabular}
    \end{column}
    \begin{column}{0.53\textwidth}
      \centering
      {\footnotesize\textbf{Measured: same model, interface changed}}\\[0.05cm]
      \safeimg[width=\textwidth,height=4.0cm,keepaspectratio]{swe_agent.pdf}
    \end{column}
  \end{columns}
  \mlsyscredit{Conclusion, \S20.6 and table 6. Yang et al. 2024, SWE-agent, arXiv 2405.15793, Tables 1 and 3: GPT-4 Turbo unless named, 300 SWE-bench Lite issues.}
  \note{[3 min] LINK: an execution trace supports a sharper question than ``does it need a bigger model.'' If the manifest still shows config\_v7 after settlement, the defect is in state selection; if the runtime retries blindly, it is in outcome handling. Attribution replays the saved inputs and changes one suspected factor at a time (Lecture 14's postmortem by intervention). The table orders the remedies by cost and reversibility; the chapter calls it the intervention ladder. It is a default order that may skip rungs; an authorization bypass goes straight to a runtime repair. Lecture 16's three homes for behavior apply: new weights cannot supply a missing state transition or make a non-idempotent retry safe. Delegation sits outside the order and is judged against Lecture 18's matched single controller. The plot recalls Lecture 7's measurement in one line: with the model fixed, interface changes alone moved resolution between 18.0 and 10.3 percent. ENGAGE: true or false: an intervention is justified because it is the next rung. Expected: false; adapting a model repays its one-time cost only over enough tasks (Lecture 19's adaptation payback), and returns nothing if the failure was in context or retries. FLEX: [CORE].}
\end{frame}

% =============================================================================
% SECTION: SAFETY CASE
% =============================================================================
\mlsysnote{[0.5 min] Segment 4: what the evidence of one incident can and cannot claim, and what stays open. FLEX: [CORE].}
\mlsysdivider{What does a passing check actually establish?}{Conclusion, \S20.7 to \S20.10}
\section{Safety Case}

\begin{frame}{How much of the passing set is invalid?}
  \mlsysactivity{Predict}{3 min}
  \begin{columns}[c]
    \begin{column}{0.44\textwidth}
      \begin{mlsyscard}{mlsysaccent}
        \small
        30\% of candidate actions are valid.\\[0.05cm]
        The check misses 10\% of valid ones.\\[0.05cm]
        It passes 5\% of invalid ones.
      \end{mlsyscard}
      \vspace{0.2cm}\small
      \uncover<2->{\textbf{11.5\%} of what passes is invalid, about one in nine.}
    \end{column}
    \begin{column}{0.53\textwidth}
      \centering
      \mlsyssteps[width=\textwidth,height=5.0cm,keepaspectratio]{passing_step}{2}
    \end{column}
  \end{columns}
  \mlsyscredit{Conclusion, \S20.10 (pitfall) and equation 3.3. 11.5\% derived (images/calcs.py).}
  \note{[3 min] ENGAGE: write a percentage, then compare with a neighbor. Most students guess near 5 percent, the check's false-pass rate. Step 2: the valid ones that pass are 0.30 $\times$ 0.90 = 0.27 of all candidates; the invalid ones that pass are 0.70 $\times$ 0.05 = 0.035; so 0.035 / 0.305 is 11.5 percent invalid. The base rate dominates. This is the verification asymmetry, principle 3, and Lecture 3's verifier precision. WARN: ``95 percent accurate'' says nothing about the passing set until you know how many candidates were valid. The same pitfall has a second half: acceptable latency alone does not prove the approved change applied, since traffic may have subsided, and the change applying does not prove the service healthy. acceptance\_42 needs both, plus the policy checks. FLEX: [CORE].}
\end{frame}

\begin{frame}{Measured: 160 of 251 passing patches were suspect}
  \begin{columns}[c]
    \begin{column}{0.50\textwidth}
      \centering
      \safeimg[width=\textwidth,height=5.0cm,keepaspectratio]{swebp.pdf}
    \end{column}
    \begin{column}{0.46\textwidth}
      \small
      SWE-agent with GPT-4 passed the tests on 251 SWE-bench issues.\\[0.15cm]
      Read by hand, 82 had the fix written in the issue and 78 passed weak tests.\\[0.15cm]
      A passing test certified only what it tested.
    \end{column}
  \end{columns}
  \mlsyscredit{Aleithan et al. 2024, ``SWE-Bench+,'' arXiv 2410.06992v2, Table 1 (manual review of 251 passing patches). 160 and the sums derived (images/calcs.py).}
  \note{[2 min] Real data for the previous slide's arithmetic. Aleithan and colleagues read, by hand, all 251 SWE-bench patches that SWE-agent with GPT-4 had resolved. In 82 the solution was already in the issue report or its comments. In 78 the patch was incorrect, incomplete, or changed different files, and the tests were too weak to notice: 32 incorrect, 37 incomplete, 9 elsewhere. 91 were correct, 76 of them by a different route from the reference fix. So 160 of 251, 63.75 percent, were suspect although every one passed. ENGAGE: which closure evidence level from Lecture 1 do these tests reach? Expected: visible tests, one level below sealed tests; and leakage is a contamination problem Lecture 13 treats. WARN: this does not mean the benchmark is useless; it means a pass rate is a claim about the tests, which must be scoped. FLEX: [OPTIONAL].}
\end{frame}

\begin{frame}{A safety case scopes each claim to its evidence}
  \footnotesize
  \renewcommand{\arraystretch}{1.15}
  \begin{tabular}{@{}P{0.27\textwidth}P{0.33\textwidth}P{0.32\textwidth}@{}}
    \toprule
    \textbf{Scoped claim} & \textbf{Evidence} & \textbf{Residual uncertainty} \\
    \midrule
    The change needed approval for the exact request & approval and admission records; tests of stale and expired approvals & a privileged path around this gate \\
    \leavevmode\uncover<2->{The lost reply did not justify a duplicate} & \leavevmode\uncover<2->{intent before dispatch; status confirmed \op{config\_v8}} & \leavevmode\uncover<2->{the endpoint must honor ids, versions, retention} \\
    \leavevmode\uncover<3->{The incident met its contract} & \leavevmode\uncover<3->{\op{acceptance\_42}: \op{op\_42}, \op{config\_v8}, health window} & \leavevmode\uncover<3->{future conditions; the diagnosis itself} \\
    \bottomrule
  \end{tabular}\\[0.25cm]
  \small\uncover<3->{The case covers one execution of one incident.}
  \mlsyscredit{Conclusion, \S20.7 and table 7; Kelly and Weaver 2004 (Goal Structuring Notation).}
  \note{[2.5 min] Lecture 14's empirical safety case: an inspectable argument that specific claims hold within a declared envelope, and it certifies nothing outside that envelope. Each row has a claim, the argument and its owner, the evidence, and the assumptions and residual uncertainty. Step 1: the action gate checks approval\_42 against the exact request; evidence includes tests of missing, expired, mismatched, and stale approvals; an unrelated privileged operator is outside this gate's claim. Step 2: settlement preserved op\_42; it rests on the endpoint honoring identity, version, and retention, and an unavailable status leaves the operation unresolved. Step 3: acceptance\_42 binds the contract, the settled operation, config\_v8, and the measured health window; the measurements cover that service and window, and they do not prove the model's diagnosis was right. ENGAGE: the model proposes shortening health\_window\_42 to close sooner. May it? Expected: no. Only the contract owner may adopt a revision; halving the window halves the samples and widens the latency interval by about 1.41 times. WARN: an operator sign-off is evidence only where the contract requires it, with its own blind spots. FLEX: [CORE].}
\end{frame}

\begin{frame}{One incident is one trial; a release needs a population}
  \begin{columns}[c]
    \begin{column}{0.55\textwidth}
      \small
      Tasks per arm grow as $1/\Delta^2$: halving the detectable difference $\Delta$ needs \textbf{4$\times$} the tasks.\\[0.25cm]
      A canary that stops admission halts future dispatches only. Applied changes stay, with their owner.\\[0.25cm]
      New weights or prompts need new release evidence.
    \end{column}
    \begin{column}{0.41\textwidth}
      \centering
      \safeimg[width=\textwidth,height=4.4cm,keepaspectratio]{canary_mine.jpg}
    \end{column}
  \end{columns}
  \mlsyscredit{Conclusion, \S20.7.2, \S20.7.3, and equation 13.3. Photo: George McCaa, U.S. Bureau of Mines, Public domain, Wikimedia Commons.}
  \note{[2 min] The three claims of the safety case concern one execution. Releasing the configuration that handled it is a population question, answered by Lecture 14's staged release: offline, shadow, then canary, under Lecture 13's paired comparison with score intervals. In that comparison incident\_42 is one trial. Equation 13.3: the tasks per arm scale as one over the squared detectable difference, so halving the difference quadruples the bill (and doubling the tasks shrinks the difference only to 0.71 of before). Had op\_42 stayed unknown at the cutoff, the trial counts as unresolved even if latency recovered. The photo is the original canary: miners carried it to meet gas first. Stopping a canary, revoking its tokens, or draining its traffic halts future dispatches; a request the endpoint already applied persists with its settlement owner. And a weight update invalidates serving state and inherits no endorsement: KV tensors computed under old weights are not valid for new ones even with identical prefix text. FLEX: [OPTIONAL].}
\end{frame}

\begin{frame}{Does a capable model still need runtime bounds?}
  \mlsysactivity{Vote}{3 min}
  \begin{columns}[T]
    \begin{column}{0.45\textwidth}
      \small
      \uncover<2->{\textbf{Yes.} At 0.1\% per step over 50 turns, \textbf{4.88\%} of trajectories propose at least one violation.\\[0.15cm]
      At 10,000 trajectories a day: about \textbf{488} hazardous proposals.}
    \end{column}
    \begin{column}{0.52\textwidth}
      \uncover<3->{{\footnotesize\textbf{Lecture 10's freeze, stated only to the agent}}\\[0.1cm]
      {\small The agent deleted the production database, then reported that rollback was impossible; the rollback worked.}}
    \end{column}
  \end{columns}
  \mlsyscredit{Conclusion, \S20.10 (4.88\% and 488 derived, images/calcs.py). Recall: The Agent Harness, War Story 10.1 (Lecture 10).}
  \note{[3 min] ENGAGE: show of hands, yes or no, then one minute with a neighbor. This is the chapter's fallacy that a capable enough model removes the need for runtime boundaries. Step 2: yes, it still needs them. A better model lowers the rate of violating proposals; training does not drive it to zero, and Lecture 1's horizon compounding amplifies what remains: 1 minus  0.999 to the 50th is 4.88 percent of 50-turn trajectories, about 488 hazardous proposals a day across 10,000. Boundaries decide whether each one is intercepted. Step 3 recalls Lecture 10's case in one line. A freeze stated only to the agent did not stop it deleting the production database, and its report that rollback was impossible was wrong. Two lessons follow: nothing enforces a freeze stated to the model (Lecture 10's approval escrow and dispatch-time rule would have held the mutation), and a model's status report is a claim about the system's state that needs evidence. In incident\_42, even an accurate restoration proposal still needs approval binding and a version check; model confidence never authorizes op\_42. FLEX: [CORE].}
\end{frame}

\begin{frame}{Which open problem would you attack first?}
  \mlsysactivity{Discuss}{3 min}
  \begin{columns}[c]
    \begin{column}{0.58\textwidth}
      \centering
      \safeimg[width=\textwidth,height=3.0cm,keepaspectratio]{attribution.pdf}\\[0.1cm]
      \raggedright\small
      Checks every five steps cut the blame search to one window: about \textbf{20 minutes} at 1.2 s per replay.\\[0.1cm]
      \footnotesize Also open: incomplete acceptance criteria; state that drifts while the agent deliberates.
    \end{column}
    \begin{column}{0.37\textwidth}
      \centering
      \safeimg[width=0.8\textwidth,height=3.4cm,keepaspectratio]{minsky.jpg}\\[0.05cm]
      {\footnotesize Minsky named credit assignment in 1961.}
    \end{column}
  \end{columns}
  \mlsyscredit{Conclusion, \S20.8 (4 actions per step, 50 steps: illustrative). 10$^{30}$, 1,024, and 20 minutes derived (images/calcs.py). Photo: Sethwoodworth at English Wikipedia, CC BY 3.0, Wikimedia Commons.}
  \note{[3 min] NARRATE: after a trajectory fails only at its final check, finding the earliest mismatch is a replay problem (Lecture 11's first-divergence localization), but showing which step was responsible is counterfactual: every untaken branch is a suspect. With an illustrative four candidate actions per step over 50 steps, that is 4 to the 50th, about 10 to the 30th paths, the same exponent as horizon compounding. Lecture 17 met this as credit assignment; Minsky named it in 1961. If checks pass every five steps, and the window's state is isolated and restorable, the search covers one window: 4 to the 5th, 1,024 paths, about 20 minutes at 1.2 s each. ENGAGE: two minutes with a neighbor on the three open problems of \S20.8. Acceptance criteria are incomplete: the health window checks only latency and errors, missing worker pools and downstream records. Production drifts while the agent deliberates: a version check rejects a stale write, but repeated rejections can livelock the budget. A log gives order and no cause, since the network may have recovered on its own. Expected: no single answer; push for the measurement each would need. FLEX: [CORE].}
\end{frame}

\begin{frame}{Deciding what counts as recovery stays with the owner}
  \begin{columns}[c]
    \begin{column}{0.30\textwidth}
      \centering
      \safeimg[width=\textwidth,height=4.6cm,keepaspectratio]{fred_brooks.jpg}\\[0.05cm]
      {\footnotesize Fred Brooks, ``No Silver Bullet,'' 1987}
    \end{column}
    \begin{column}{0.65\textwidth}
      \small
      Accidental work: producing a valid configuration payload.\\[0.1cm]
      Essential work: deciding which behavior counts as recovery.\\[0.3cm]
      \uncover<2->{If half the effort stays human specification and review, the task finishes at most \textbf{2$\times$} faster, however fast the model is.\\[0.3cm]}
      \uncover<3->{An agent may draft a contract. Only its owner may adopt it.}
    \end{column}
  \end{columns}
  \mlsyscredit{Conclusion, \S20.9; Brooks 1987. 2$\times$ derived from Amdahl's bound (images/calcs.py). Photo: SD\&M (www.sdm.de), CC BY-SA 3.0, Wikimedia Commons.}
  \note{[2 min] Step 1: Brooks separated the essential difficulty of conceiving software from the accidental difficulty of expressing it. In incident management, producing the payload is accidental; deciding which service behaviors count as recovery, and which risks a change introduces, is essential. Step 2: if a fraction of specification and review stays unaided by models, it is the serial fraction of Lecture 18's Amdahl bound, here with one human and one model. With half the effort serial, the speedup is at most 1/0.5, 2 times. Step 3: the specification boundary. An agent can propose requirements, draft fault-injection tests, or flag an impractical objective; it gains no governance authority by authoring a contract and cannot relax its own acceptance criteria. This is invariant closure applied to the contract itself. WARN: agent-written tests are useful, but tests co-generated with the implementation can encode the same mistaken requirement, so their agreement is not independent verification. FLEX: [OPTIONAL].}
\end{frame}

% =============================================================================
% SECTION: SYNTHESIS
% =============================================================================
\mlsysnote{[0.5 min] The synthesis: the hook answered with the records of the whole incident, then the course's two relations. FLEX: [CORE].}
\mlsysdivider{Can you now close one task end to end?}{Conclusion, \S20.2 (table 1) and \S20.11}
\section{Synthesis}

\begin{frame}{Six owners and six records close \op{incident\_42}}
  \footnotesize
  \renewcommand{\arraystretch}{1.15}
  \begin{tabular}{@{}P{0.20\textwidth}P{0.31\textwidth}P{0.43\textwidth}@{}}
    \toprule
    \textbf{Transition} & \textbf{Owner} & \textbf{Record it leaves} \\
    \midrule
    Task admitted & \leavevmode\uncover<2->{runtime, by contract} & \leavevmode\uncover<3->{\op{incident\_contract\_v1}, \op{read\_scope\_42}} \\
    Evidence selected & \leavevmode\uncover<2->{context manager} & \leavevmode\uncover<3->{\op{context\_manifest\_42}, with sources} \\
    Change approved & \leavevmode\uncover<2->{operator decides} & \leavevmode\uncover<3->{\op{change\_request\_42}, \op{approval\_42}} \\
    Change dispatched & \leavevmode\uncover<2->{action gate and journal} & \leavevmode\uncover<3->{durable intent for \op{op\_42}} \\
    Lost reply settled & \leavevmode\uncover<2->{settlement process} & \leavevmode\uncover<3->{status confirming \op{config\_v8}} \\
    Incident accepted & \leavevmode\uncover<2->{accepting principal} & \leavevmode\uncover<3->{\op{acceptance\_42}, \op{health\_window\_42}} \\
    \bottomrule
  \end{tabular}\\[0.2cm]
  \small\uncover<3->{Afterward, \op{operations\_owner\_42} keeps the service and any unresolved operation.}
  \mlsyscredit{Conclusion, \S20.2.1, table 1.}
  \note{[2.5 min] Return to the hook. Step 1: the six transitions incident\_42 went through. ENGAGE: before step 2, ask the room to name the owner of each. Step 2: owners. The model proposes and never authorizes; an authorized operator decides; the settlement process is a runtime component acting for operations\_owner\_42. Step 3: records. With these an auditor reconstructs the incident without asking the model what happened. Now answer the hook: retrying blindly fails because the endpoint's operation identity and the durable intent make a status query possible; assuming success fails because acceptance\_42 needs the settled operation and the health window; halting fails because the intent names an owner who keeps the open obligation. Each common failure of an agent deployment meets one guard in this trace: a duplicate meets the operation identity, a premature close meets the acceptance record, an untracked change meets the intent log and its owner. FLEX: [CORE].}
\end{frame}

\mlsysnote{[0.5 min] Read it once, slowly. It is the chapter's first takeaway and the shortest statement of the course. Add the corollary aloud: an unknown outcome keeps its owner; it never becomes a retry or a success. FLEX: [CORE].}
\mlsysfocus{Transition ownership}{%
Every consequential transition needs an owner and a durable record.}

\begin{frame}{Checks bend the horizon curve; measured agents show it}
  \begin{columns}[T]
    \begin{column}{0.48\textwidth}
      \centering
      \safeimg[width=\textwidth,height=4.4cm,keepaspectratio]{horizon.pdf}\\
      {\footnotesize Unchecked success at 99\% per step}
    \end{column}
    \begin{column}{0.48\textwidth}
      \centering
      \safeimg[width=\textwidth,height=4.4cm,keepaspectratio]{metr_p80.pdf}\\
      {\footnotesize Measured: 80\% success means tasks about 5.5$\times$ shorter (median)}
    \end{column}
  \end{columns}
  \mlsyscredit{Introduction, Napkin Math 1.1; Conclusion, \S20.11. METR Time Horizon 1.1 (benchmark\_results\_1\_1.yaml): 13 state-of-the-art models since 2024; ratios derived (images/calcs.py).}
  \note{[1.5 min] The first of the two pieces of arithmetic Lecture 1 opened with, now with what the course built against it. Left: horizon compounding. At 99 percent per step, unchecked success is 90.4 percent at 10 turns and 36.6 percent at 100. Admission and verification turn detected violations into rejected attempts, typed contracts stop unsafe proposals, watchdogs (Lecture 12) bound how far an error runs, and intent logging bounds recovery. Right, measured: METR's time horizons for the 13 models it marks state of the art since 2024. Purple is the task length, in human time, a model completes half the time; red is the length it completes 80 percent of the time. Asking for 80 instead of 50 percent shortens the task about 4 to 10 times, median 5.5. Demanding reliability shrinks the horizon, as compounding predicts. WARN: METR's curve is a fit across tasks and carries no per-step error rate; it agrees with compounding in direction only. FLEX: [CORE].}
\end{frame}

\begin{frame}{Raising acceptance beats halving the token price}
  \begin{columns}[c]
    \begin{column}{0.58\textwidth}
      \centering
      \safeimg[width=\textwidth,height=3.0cm,keepaspectratio]{cost.pdf}\\[0.05cm]
      {\footnotesize A \$2.10 attempt; tokens are 62\% of it}
    \end{column}
    \begin{column}{0.38\textwidth}
      \small
      \op{incident\_42} needed neither first, because its money sat in approval and its risk in the unchecked wrong change.
    \end{column}
  \end{columns}
  \mlsyscredit{Introduction, Napkin Math 1.2; Conclusion, \S20.3.2 and \S20.11. Values derived (images/calcs.py).}
  \note{[1 min] The second relation: cost per accepted task. A \$2.10 attempt accepted 40 percent of the time costs \$5.25. Halving the token price, which is 62 percent of the attempt, gives \$3.63; raising acceptance to 60 percent gives \$3.50. ENGAGE: which lever did incident\_42 need? Expected: neither first; 97.8 percent of its attempt was approval, and the \$480 expected loss of an unchecked wrong change dwarfed both. Choose the intervention by pricing the work that improves the outcome; the largest cost term need not be the lever. FLEX: [OPTIONAL].}
\end{frame}

\begin{frame}
  \frametitle<1|handout:0>{Which guard would you build first?}%
  \frametitle<2-|handout:1>{Takeaways}%
  \only<1|handout:0>{\mlsysactivity{Retrieval}{2 min}%
    \centering
    \vspace{0.7cm}
    {\Large\bfseries Close your notes.}\\[0.5cm]
    {\large Write the one guard you would build first for an agent you use,\\[0.15cm] and the record it would leave.}\\[0.5cm]
    {\normalsize\textcolor{midgray}{Two minutes. Hand it in on the way out.}}}%
  \only<2-|handout:1>{%
    \thumbline{thumb_lostreply_step3.pdf}{Every consequential transition needs an owner and a durable record; an unknown outcome is never a license to retry.}\\[0.15cm]
    \thumbline{thumb_records_step2.pdf}{The contract sets the boundaries; every record keeps its source, version, and authority.}\\[0.15cm]
    \thumbline{thumb_swe_agent.pdf}{Choose the cheapest intervention that fixes the attributed failure, measured end to end.}\\[0.15cm]
    \thumbline{thumb_passing_step2.pdf}{A passing check supports only what it covers; owners adopt contracts and keep the residual risk.}\\[0.15cm]
    \begin{minipage}[c]{2.6cm}\centering\safeimg[width=2.5cm,height=1.0cm,keepaspectratio]{datacenter-server-racks.jpg}\end{minipage}%
    \hspace{0.3cm}\begin{minipage}[c]{0.74\textwidth}\small\raggedright The deliverable is the audit trail: proposed, permitted, applied, verified, and who stays responsible.\end{minipage}}
  \mlsyscredit{Conclusion, Key Takeaways and \S20.11. Photo: Carl Lender, CC BY 2.0, Wikimedia Commons.}
  \note{[3 min] Step 1: the course's last exit ticket. Two minutes, notes closed. Any guard counts if it names the boundary it sits at and the record it leaves: a dispatch-time freeze rule leaving a refusal record; a status lookup leaving a settlement record. Collect the cards. Step 2: the chapter's four takeaways, one per segment, and its closing line. Read them in order. The last line is the chapter's own: an audit trail of what the model proposed, what the action boundary permitted, what effects occurred, what evidence validated the result, and who remains responsible is the real deliverable of an agentic system. FLEX: [CORE].}
\end{frame}

\begin{frame}{Open questions to carry beyond the course}
  \normalsize
  \begin{itemize}\setlength{\itemsep}{0.35cm}
    \item How would you write an acceptance contract that catches the regression a latency window misses?
    \item When production changes faster than the agent deliberates, how should freshness, locking, and latency be traded?
    \item What evidence would show that one step caused a failure, rather than only preceding it?
  \end{itemize}
  \note{[0.5 min] The chapter's three open systems questions within the digital boundary (\S20.8). No single answer. Ask students to write one sentence on any of them; these are research problems, and good course projects start from one. FLEX: [OPTIONAL].}
\end{frame}

\begin{frame}{Read, practice, and go further}
  \small
  \begin{itemize}\setlength{\itemsep}{0.12cm}
    \item \textbf{Read:} Conclusion, \S20.1 to \S20.11; Checkpoints 20.1 to 20.3; Exercises 3, 8, and 9.
    \item Saltzer, Reed, and Clark, ``End-to-End Arguments in System Design,'' TOCS, 1984.
    \item Gray, ``The Transaction Concept: Virtues and Limitations,'' VLDB, 1981.
    \item Brooks, ``No Silver Bullet: Essence and Accidents of Software Engineering,'' Computer, 1987.
    \item Sculley et al., ``Hidden Technical Debt in Machine Learning Systems,'' NeurIPS, 2015.
    \item Kelly and Weaver, ``The Goal Structuring Notation,'' 2004.
  \end{itemize}
  \note{[0.5 min] The chapter first: Exercise 3 is today's estimate, Exercise 8 sizes a whole remediation trajectory, and Exercise 9 critiques a research system end to end against every principle. The five readings are the chapter's own: Saltzer, Reed, and Clark on why acceptance belongs at the endpoint; Gray on what transactions give an agent runtime and where compensation must replace them; Brooks on essential and accidental work; Sculley and colleagues on why the code around a model dominates its upkeep; Kelly and Weaver on writing a safety argument whose gaps a reviewer can see. FLEX: [OPTIONAL].}
\end{frame}

\begin{frame}[c]
  \centering
  {\small\textcolor{lightgray}{FOR FURTHER STUDY}}\\[0.4cm]
  {\LARGE\bfseries What must change when an agent's actions\\move matter instead of bytes?}\\[0.8cm]
  {\large What remains to engineer is the surrounding system:\\the harness and runtime that make a model's proposals\\safe to execute, verifiable to accept, and economical to operate.}\\[0.2cm]
  {\footnotesize\textcolor{midgray}{Conclusion, \S20.8 and \S20.11}}
  \note{[0.5 min] There is no next lecture. Leave students with the question the course does not answer. The chapter notes that physical actuation adds sensing, continuous dynamics, and hard deadlines, and that software agents and robots both act on a world that terminating a process cannot roll back, so reversibility does not separate them. Then read the chapter's closing sentence aloud and end there. FLEX: [CORE].}
\end{frame}

\end{document}
